% =============================================================================
%  METHODS.tex
%  Detailed methodology of the protein-perturbation virtual-screening pipeline.
%
%  This file is written as plain LaTeX so it can be compiled with `pdflatex`
%  but is also fully readable in a text editor. Keep it next to README.md.
%
%  Compile:  pdflatex METHODS.tex
% =============================================================================

\documentclass[11pt]{article}

\usepackage[a4paper, margin=2.5cm]{geometry}
\usepackage{amsmath, amssymb, amsthm}
\usepackage{mathtools}
\usepackage{hyperref}
\usepackage{booktabs}
\usepackage{listings}
\usepackage{xcolor}
\usepackage{enumitem}

\lstset{
    basicstyle=\ttfamily\small,
    columns=fullflexible,
    breaklines=true,
    keepspaces=true,
    backgroundcolor=\color{gray!10},
    frame=single,
    framesep=4pt,
    rulecolor=\color{gray!40},
}

\newcommand{\jsd}{\mathrm{JSD}}
\newcommand{\kl}{D_{\mathrm{KL}}}
\newcommand{\prob}{\mathbb{P}}
\newcommand{\E}{\mathbb{E}}
\newcommand{\Indic}{\mathbf{1}}

\title{The Protein-Perturbation Score \\
       \large A single-shot, MD-free virtual-screening method \\
       based on LigandMPNN inverse-folding distributions}
\author{Guglielmo Tedeschi}
\date{\today}

\begin{document}
\maketitle

\begin{abstract}
We assess whether the change a ligand induces in the per-residue amino-acid
distribution produced by LigandMPNN, relative to a co-evolved reference
complex, can be used as a discrimination signal between active and decoy
binders. The method requires only a single forward pass of LigandMPNN per
docked pose (no molecular dynamics, no free-energy perturbation, no
finetuning), and reduces every pose to a single scalar: the mean
Jensen--Shannon divergence across the pocket, which we call the
\emph{perturbation score}. We benchmark the score on an ABL1 / imatinib
active--decoy set, and report AUROC, enrichment factors, and a bootstrap
confidence interval. The framework is target-agnostic: switching protein
requires only changing the reference complex and the docked library.
\end{abstract}

\section{Motivation}

Structure-based virtual screening typically falls into two regimes:
\emph{fast} (docking scores, shape complementarity) and \emph{slow} (free
energy of binding via MD/FEP). The fast methods are inexpensive but noisy;
the slow methods are accurate but inapplicable at library scale.

LigandMPNN~\cite{LigandMPNN} is an inverse-folding model: given a backbone
and a ligand, it outputs, for each residue position, a probability
distribution over the 20 canonical amino acids. The distribution reflects
which residue identity the network believes is most consistent with the
\emph{backbone + ligand} context. A pocket that has co-evolved with its
native ligand yields, at the pocket residues, a low-entropy distribution
peaked on the wild-type residue. A non-native ligand placed in the same
pocket pushes that distribution away from its peak.

The central hypothesis of this work is:

\begin{quote}
\emph{The magnitude of the per-residue distribution drift, measured between
the reference complex and a candidate pose, is small for active binders
and large for decoys.}
\end{quote}

If true, the drift itself becomes a usable ranking signal --- one that is
free of explicit binding-energy estimation, requires no sampling, and is
computed in a single LigandMPNN forward pass.

\section{Notation}

\begin{tabular}{ll}
\toprule
Symbol & Meaning \\
\midrule
$N$ & Number of residues in the pocket selection (identical for ref \& screening)\\
$A=21$ & Alphabet size (20 canonical AAs + unknown token \texttt{X})\\
$R$ & Number of LigandMPNN replicas per pose \\
$P_{\mathrm{ref}}\in [0,1]^{N\times A}$ & Reference per-residue AA distribution \\
$Q_p\in [0,1]^{N\times A}$ & Distribution for screening pose $p$ \\
$\jsd_i$ & Jensen--Shannon divergence at residue $i$ \\
$S_p$ & Perturbation score for pose $p$ \\
$y_p \in \{0,1\}$ & Label: $1$ = active, $0$ = decoy \\
\bottomrule
\end{tabular}

All probability matrices are row-stochastic:
$\sum_{a=1}^{A} P_{\mathrm{ref}}[i,a] = 1$ for every residue $i$.

\section{Pipeline}

The pipeline is a linear sequence of six stages.

\begin{enumerate}[leftmargin=*, label=\textbf{(\arabic*)}]

  \item \textbf{Docking (\texttt{gnina\_docking.py}).}
  Each ligand of the library is docked into the protein with
  GNINA~\cite{GNINA}, using the reference ligand to define the binding
  site:
  \[
  \mathrm{box\,center} = \frac{1}{N_{\mathrm{lig}}}\sum_{k=1}^{N_{\mathrm{lig}}}
  \mathbf{r}_k^{\mathrm{ref}},
  \qquad
  \mathrm{box\,size} = \mathrm{ref\,extent} + 2\,\Delta_{\mathrm{box}}.
  \]
  $\Delta_{\mathrm{box}}$ (\texttt{autobox\_add}, default 8\,\AA) is the
  isotropic padding around the reference. GNINA returns up to
  $\texttt{num\_modes}=9$ poses per ligand, each annotated with three
  scores: \texttt{CNNscore}, \texttt{CNNaffinity}, and
  \texttt{minimizedAffinity}.

  \item \textbf{Pose selection (\texttt{combine\_sdf\_protein.py}).}
  For each ligand, we keep the single pose with the highest
  \texttt{CNNaffinity} (or another GNINA-reported metric, configurable)
  and write a protein--ligand complex PDB.

  \item \textbf{Pose QC (\texttt{pose\_qc.py}).}
  Two checks are run on every complex:

  \begin{itemize}[leftmargin=*]
    \item \emph{Clash detection (FAIL).} For every pair (protein heavy
    atom $a$, ligand heavy atom $b$), we compute the distance $d_{ab}$ and
    compare it to an element-aware threshold
    \[
    \tau_{ab} = \max\bigl(\tau_{\min},\; 0.75\,(r^{\mathrm{cov}}_a +
    r^{\mathrm{cov}}_b)\bigr),
    \]
    where $r^{\mathrm{cov}}$ are tabulated covalent radii. The pose fails
    if any $d_{ab} < \tau_{ab}$. The $0.75$ factor and the
    \texttt{clash\_dist} floor ($\tau_{\min}=1.5$\,\AA\ by default)
    together prevent both false negatives (true overlaps) and false
    positives (genuine covalent linkers).
    \item \emph{Geometry check (WARN).} We rebuild an RDKit molecule from
    the ligand PDB block and verify, per bond, that the length lies in a
    tabulated $(\ell_{\min},\ell_{\max})$ range, and per heavy-atom triple
    that the angle lies in a tabulated $(\theta_{\min},\theta_{\max})$
    range. Violations raise warnings; under \texttt{--strict} they become
    failures.
  \end{itemize}

  Failed poses are dropped before LigandMPNN sees them.

  \item \textbf{Pocket residue selection (\texttt{residues\_selection.py}).}
  Given the reference complex, we compute the geometric center
  $\mathbf{c}$ of the reference ligand and select every protein residue
  with at least one atom inside the autobox sphere:
  \[
  \mathcal{R} \;=\; \bigl\{ \text{residue $i$} \;:\;
    \min_{j \in \text{atoms}(i)} \|\mathbf{r}_j - \mathbf{c}\|
    \le \Delta_{\mathrm{box}} \bigr\}.
  \]
  $\mathcal{R}$ is fixed across the whole library, so reference and
  screening matrices have identical shape and the per-residue comparison
  is well-defined.

  \item \textbf{LigandMPNN forward pass (\texttt{run\_pipeline.sh},
  \texttt{run\_score.sh}).} For each complex $p$, LigandMPNN is run with
  $R$ replicas (different seeds, identical inputs) over the residues
  $\mathcal{R}$, producing a tensor of shape $(R, N, A)$. We average over
  replicas to reduce stochastic noise:
  \[
  Q_p[i, a] \;=\; \frac{1}{R} \sum_{r=1}^{R} \mathrm{LigandMPNN}_r(p)[i, a].
  \]
  The reference complex is run identically, producing
  $P_{\mathrm{ref}}$.

  \item \textbf{Perturbation score \& evaluation (\texttt{aucroc.py}).}
  See Section~\ref{sec:score}.

\end{enumerate}

\section{The perturbation score}
\label{sec:score}

\subsection{Per-residue divergence}

We use the Jensen--Shannon divergence, base 2:
\begin{equation}
\jsd(P \,\|\, Q) \;=\; \tfrac{1}{2}\,\kl(P \,\|\, M)
                  + \tfrac{1}{2}\,\kl(Q \,\|\, M),
\qquad
M = \tfrac{1}{2}(P + Q),
\label{eq:jsd}
\end{equation}
with $\kl(P\|Q) = \sum_a P(a)\log_2\!\bigl(P(a)/Q(a)\bigr)$.

JSD has three properties that matter here:
\begin{itemize}[leftmargin=*]
  \item \emph{Symmetric}, $\jsd(P,Q)=\jsd(Q,P)$, so the score does not
        depend on which side we call "reference".
  \item \emph{Bounded} in $[0, 1]$ when $\log_2$ is used. The score is
        thus comparable across targets and residues without rescaling.
  \item \emph{Smoothing-free.} Since $M=\tfrac12(P+Q)$ is strictly
        positive wherever either $P$ or $Q$ is positive, no
        $\varepsilon$ regularisation is needed --- unlike a naive
        $\kl(P\|Q)$ which diverges if $Q$ assigns zero probability to a
        residue identity supported by $P$.
\end{itemize}

\subsection{Aggregated pose score}

Given the reference $P_{\mathrm{ref}}$ and pose distribution $Q_p$ on the
same $N$ residues, we compute a JSD vector
$\bigl(\jsd_1^{(p)},\dots,\jsd_N^{(p)}\bigr)$ via Eq.~\eqref{eq:jsd}, and
aggregate to a single scalar:
\begin{equation}
S_p \;=\; 100 \times \frac{1}{N} \sum_{i=1}^{N} \jsd_i^{(p)}.
\label{eq:score}
\end{equation}
The factor 100 is cosmetic --- it puts the score on a $0$--$100$ scale
since $\jsd_i \in [0,1]$.

\subsection{Sign of the score}

By construction, $S_p \ge 0$. The screening hypothesis is that active
ligands give small $S_p$ and decoys give large $S_p$. To use $S_p$
directly in a ROC curve, where higher values must indicate the positive
class, we work with $-S_p$:
\begin{equation}
\hat s_p \;\equiv\; -S_p, \qquad
y_p \;=\; \begin{cases} 1 & \text{if pose $p$ is an active}, \\ 0 & \text{if a decoy}. \end{cases}
\end{equation}

\subsection{Duplicate compounds}

A library entry may produce several poses (e.g.\ stereoisomers,
tautomers). We keep, per compound id, the pose with the \emph{lowest}
$S_p$ (the most active-like) and discard the others. This is the
standard "best-pose-per-compound" convention; it makes the AUROC reflect
ranking power at the compound level, not at the pose level.

\section{Statistical evaluation}

\subsection{AUROC}

The Area Under the Receiver-Operating-Characteristic curve is computed
on $(\hat s_p, y_p)$ pairs in the usual way:
\begin{equation}
\mathrm{AUROC} \;=\;
\frac{1}{n_+\,n_-}
\sum_{p : y_p=1}\;\sum_{q : y_q=0}
\Indic\!\bigl[\hat s_p > \hat s_q\bigr] +
\tfrac12\,\Indic\!\bigl[\hat s_p = \hat s_q\bigr],
\end{equation}
where $n_+$ ($n_-$) is the number of actives (decoys). It is the
probability that, drawing one active and one decoy at random, the active
has the higher predicted score.

\subsection{Bootstrap confidence interval}

We resample with replacement the $(y_p, \hat s_p)$ pairs $B$ times
(default $B=2000$, seed $42$); on each resample we recompute the AUROC,
skipping degenerate resamples (only one class). The percentile bootstrap
$95\%$ CI is
\[
\bigl[\mathrm{AUROC}_{(2.5)},\;\mathrm{AUROC}_{(97.5)}\bigr],
\]
i.e.\ the empirical $2.5$ and $97.5$ percentiles of the $B$ resampled
AUROCs.

\subsection{Enrichment factor}

For a fraction $f \in (0,1]$ (default $f \in \{0.01, 0.05, 0.10\}$):
\begin{equation}
\mathrm{EF}(f) \;=\;
\dfrac{
    \#\{p : y_p=1 \text{ and } \hat s_p \in \text{top-}f \}
    \;/\;
    \#\{\text{top-}f\}
}{
    n_+ / n
},
\qquad
\#\{\text{top-}f\} = \lfloor f \cdot n \rfloor.
\end{equation}
$\mathrm{EF}(f)=1$ means "no better than random"; the theoretical maximum
is $\min(1, n_+/\#\{\text{top-}f\}) / (n_+/n)$ and depends on the library
composition.

\subsection{Operating point: Youden's J}

Picking a single threshold turns the ranker into a classifier. We use
Youden's index, the threshold $t^*$ that maximises
\[
J(t) \;=\; \mathrm{TPR}(t) - \mathrm{FPR}(t).
\]
At that threshold we report accuracy, balanced accuracy, precision,
recall, specificity, $F_1$, and Matthews' correlation coefficient (MCC).

\section{Implementation map}

\begin{center}
\begin{tabular}{ll}
\toprule
Stage & Script \\
\midrule
Docking                        & \texttt{scripts/gnina\_docking.py} \\
Best-pose selection            & \texttt{scripts/combine\_sdf\_protein.py} \\
Clash + geometry QC            & \texttt{scripts/pose\_qc.py} \\
Pocket residue selection       & \texttt{scripts/residues\_selection.py} \\
LigandMPNN design batch        & \texttt{scripts/run\_pipeline.sh} \\
LigandMPNN per-pose scoring    & \texttt{scripts/run\_score.sh} \\
JSD score, AUROC, EF, CI       & \texttt{scripts/aucroc.py} \\
Single-pair inspection         & \texttt{analysis.ipynb} \\
Config (single source of truth)& \texttt{scripts/config.yaml} \\
YAML reader (PyYAML)           & \texttt{scripts/parse_config.py} \\
\bottomrule
\end{tabular}
\end{center}

\section{Operating assumptions and failure modes}

\paragraph{Pocket invariance.}
The same residue set $\mathcal{R}$ is used for the reference and for every
screening pose, by construction. If the user re-runs the selection with a
different \texttt{autobox\_add}, all previously computed $.pt$ files
become incomparable to the new reference --- $Q_p$ and $P_{\mathrm{ref}}$
would have different $N$. The pipeline detects shape mismatches and
skips, but a partial regeneration leaves a silently inconsistent dataset.

\paragraph{Temperature.}
The LigandMPNN sampling temperature $T$ flattens (high $T$) or sharpens
(low $T$) the per-residue distributions. JSD is not invariant under $T$:
absolute scores depend on it. AUROC is approximately stable but should
not be reported across runs with different $T$.

\paragraph{Reference choice.}
A single reference complex carries all the structural assumptions. For
targets where the apo-to-holo motion is large, using a single crystal as
reference may bias the distribution toward one conformational state.
\texttt{aucroc.py} accepts a folder as \texttt{--reference}, in which
case the per-file distributions (each averaged over its own replicas) are
averaged a second time across files --- this is the right thing if you
want to use multiple snapshots of the holo state.

\paragraph{Active/decoy ratio.}
DUD-E-style benchmarks have $\approx 50\!:\!1$ decoy:active ratios; the
ABL1 set used here is smaller and more balanced. The Enrichment Factor
ceilings reported in the figure account for the actual ratio.

\paragraph{Pose quality dominates.}
Garbage in, garbage out: a docking pose with the ligand outside the
pocket will trivially produce a small perturbation (because LigandMPNN
sees no ligand in the binding site) and will look like an active. The
QC stage (clashes + geometry) is the main guard against this; users
who disable it should expect score inflation in the active-like
direction.

\paragraph{Validation set is small.}
At the time of writing the ABL1 set contains $\mathcal{O}(10^1)$ actives;
AUROC point estimates therefore have wide bootstrap CIs, and the
qualitative direction of the effect (small $S$ = active) is the main
claim --- not the precise AUROC value.

\section{Reproducibility checklist}

\begin{itemize}[leftmargin=*]
  \item \textbf{Seeds.} LigandMPNN design: $\mathrm{seed}=42$
    (\texttt{config.yaml}); LigandMPNN scoring: $\mathrm{seed}=111$
    (\texttt{run\_score.sh}); bootstrap: $\mathrm{seed}=42$
    (\texttt{aucroc.py}).
  \item \textbf{Environments.} \texttt{environment-analysis.yml} and
    \texttt{environment-ligandmpnn.yml} pin the runtime; the LigandMPNN
    checkpoint hash should also be recorded per run.
  \item \textbf{Inputs.} Reference complex PDB (cleaned, single chain,
    single ligand), docked SDFs (one per ligand), pocket residue list
    (auto-derived from the reference).
  \item \textbf{Outputs.} For every screening pass:
    a CSV of per-pose scores, a 3-panel ROC/score/EF PNG, and the
    LigandMPNN \texttt{.pt} files (kept under \texttt{output\_*/} or
    \texttt{experiments/}, both git-ignored).
\end{itemize}

\section{Glossary}

\begin{description}[leftmargin=2.5cm, style=nextline]
  \item[Active] A ligand annotated as a real binder of the target. Label
    $y=1$.
  \item[Decoy] A ligand annotated as a non-binder, typically chosen with
    physico-chemical properties matched to the actives (DUD-E style).
    Label $y=0$.
  \item[Pocket] The set $\mathcal{R}$ of residues within the autobox
    sphere centred on the reference ligand.
  \item[Perturbation score] $S_p$ as in Eq.~\eqref{eq:score}; the
    pose-level scalar fed into the ROC.
  \item[Reference complex] The protein--ligand crystal (or otherwise
    high-confidence structure) used to define both the pocket and the
    reference distribution $P_{\mathrm{ref}}$.
  \item[Replica] One LigandMPNN forward pass with a fresh seed; replicas
    are averaged before any divergence is computed.
\end{description}

\begin{thebibliography}{9}

\bibitem{LigandMPNN}
Dauparas, J., Lee, G.R., Pecoraro, R., An, L., Anishchenko, I., Glasscock,
C., Baker, D.
\emph{Atomic context-conditioned protein sequence design using LigandMPNN.}
Nature Methods, 2025.

\bibitem{GNINA}
McNutt, A.T., Francoeur, P., Aggarwal, R., Masuda, T., Meli, R., Ragoza,
M., Sunseri, J., Koes, D.R.
\emph{GNINA 1.0: molecular docking with deep learning.}
Journal of Cheminformatics 13, 43 (2021).

\end{thebibliography}

\end{document}
